\chapter{Wielokrotne testowanie: FWER, FDR i decyzje po selekcji}
\label{ch:multiple-testing}
\paragraph{Powiązania i odsyłacze.} Przed użyciem korekt warto przeczytać rozdziały~\ref{ch:testing} i~\ref{ch:confidence-intervals}; do projektowania badań przejdziemy w rozdziale~\ref{ch:causal-ab}.


\section{Dlaczego test wykonany poprawnie może prowadzić do błędnej procedury?}
Testowanie jednej hipotezy z poziomem $\alpha=0{,}05$ jest
różne od testowania stu hipotez, z których następnie
wybieramy tylko te o małych wartościach p. Każdy test
może być poprawny oddzielnie, a cała procedura może mieć
duże prawdopodobieństwo przynajmniej jednego
fałszywego odkrycia. Trzeba więc ustalić, \emph{jaki rodzaj
błędu kontrolujemy w zbiorze decyzji}.

\begin{definicja}[Rodzina hipotez i miary błędów]\label{def:fwer-fdr}
Niech $m\ge1$, a $H_1,\ldots,H_m$ będą ustaloną rodziną
hipotez. Niech $I_0\subseteq\{1,\ldots,m\}$ oznacza
zbiór indeksów prawdziwych hipotez zerowych i
$m_0=|I_0|$ (nieznane badaczowi). Dla testów
$\varphi_i\in\{0,1\}$ definiujemy liczbę wszystkich
odrzuceń $R=\sum_{i=1}^m\varphi_i$ oraz liczbę
błędnych odrzuceń $V=\sum_{i\in I_0}\varphi_i$.
\emph{Rodzinne prawdopodobieństwo błędu I rodzaju} to
$\mathrm{FWER}=P(V\ge1)$, a \emph{oczekiwany udział
fałszywych odkryć} to
\[
 \mathrm{FDR}=\E\left[\frac{V}{R\vee1}\right],
 \qquad R\vee1:=\max\{R,1\}.
\]
Konwencja $V/(R\vee1)=0$ dla $R=0$ jest częścią definicji.
Dodatkowo $\mathrm{FDP}=V/(R\vee1)$ oznacza
\emph{zrealizowany} udział fałszywych odkryć.
\end{definicja}

FWER dotyczy zdarzenia, że \emph{jakikolwiek} wniosek
jest fałszywy. FDR dopuszcza pewną liczbę błędnych
odkryć, ale kontroluje \emph{średni} ich udział.
Nie oznacza to, że w danym konkretnym badaniu udział
błędnych odkryć jest nie większy niż $5\%$.

\begin{definicja}[Silna i słaba kontrola FWER]
Procedura ma \emph{silną kontrolę FWER} na poziomie
$\alpha$, jeśli $P(V\ge1)\le\alpha$ dla każdego
rozmieszczenia prawdziwych i fałszywych hipotez w rozważanej
rodzinie modeli. \emph{Słaba kontrola FWER} wymaga tej
nierówności tylko w konfiguracji, w której wszystkie
hipotezy zerowe są prawdziwe.
\end{definicja}

\begin{przyklad}[Dwadzieścia niezależnych testów]
Jeśli pod prawdziwymi hipotezami zerowymi
$p_1,\ldots,p_{20}$ są niezależne i jednostajne na
$[0,1]$, a każde $H_i$ odrzucamy przy $p_i\le0{,}05$,
to $\mathrm{FWER}=1-(1-0{,}05)^{20}\approx0{,}6415$.
Nawet jeśli każdy pojedynczy test ma poprawny poziom,
cała procedura nie kontroluje rodzinnego błędu na $5\%$.
Ten konkretny wzór wykorzystuje \emph{niezależność}; bez
niej nie można go automatycznie stosować.
\end{przyklad}

\section{Bonferroni: prosta gwarancja bez niezależności}
\begin{twierdzenie}[Procedura Bonferroniego]\label{thm:bonferroni}
Niech $p_1,\ldots,p_m$ będą mierzalnymi wartościami p
dla ustalonej rodziny hipotez i niech dla każdego
$i\in I_0$ oraz każdego $u\in[0,1]$ zachodzi
$P(p_i\le u)\le u$. Bez żadnych założeń o wzajemnej
zależności wartości p reguła
\[
 \text{odrzuć }H_i\quad\Longleftrightarrow\quad
 p_i\le\frac\alpha m
\]
ma silną kontrolę $\mathrm{FWER}\le\alpha$.
\end{twierdzenie}
\begin{proof}
Zdarzenie $\{V\ge1\}$ jest zawarte w unii zdarzeń
$\{p_i\le\alpha/m\}$ po $i\in I_0$.
Z nierówności sumy prawdopodobieństw i ważności
poszczególnych wartości p otrzymujemy
\[
 P(V\ge1)\le\sum_{i\in I_0}P(p_i\le\alpha/m)
 \le\frac{m_0\alpha}m\le\alpha.
\]
Dowód nie wykorzystuje niezależności ani liczby fałszywych
hipotez w rodzinie.
\end{proof}

\begin{wniosek}[Jednoczesne przedziały ufności]
\label{cor:simultaneous-ci}
Niech $\theta_1,\ldots,\theta_m$ oznaczają ustalone
parametry (lub funkcjonały parametrów), a losowe
przedziały $C_i(X)$ spełniają
$P(\theta_i\notin C_i(X))\le\alpha/m$ dla każdego
$i=1,\ldots,m$. Wówczas
\[
 P\bigl(\theta_i\in C_i(X)\ \text{dla wszystkich }
 i=1,\ldots,m\bigr)\ge1-\alpha.
\]
\end{wniosek}
\begin{proof}
Zdarzenie braku jednoczesnego pokrycia jest unią zdarzeń
$\{\theta_i\notin C_i(X)\}$. Nierówność sumy
prawdopodobieństw ogranicza jego prawdopodobieństwo przez
$\sum_i\alpha/m=\alpha$. Żadna niezależność przedziałów
nie jest potrzebna.
\end{proof}

\begin{uwaga}[Cena ochrony rodziny]
Bonferroni zachowuje gwarancję przy dowolnej zależności,
ale może być konserwatywny, zwłaszcza przy wielu silnie
skorelowanych testach. Jeżeli mamy z góry wybrany
jeden \emph{główny} wynik badania i wiele wyników
pomocniczych, rozsądny plan może oddzielać rodziny
hipotez o różnych celach; nie wolno jednak tworzyć
takich rodzin dopiero po obejrzeniu wartości p, tylko
po to, by uniknąć korekty.
\end{uwaga}

\section{Procedura Holma: lepsza od Bonferroniego bez dodatkowych założeń}
\begin{definicja}[Procedura krokowa Holma]\label{def:holm}
Niech $p_{(1)}\le\cdots\le p_{(m)}$ będą uporządkowanymi
wartościami p, a $H_{(i)}$ odpowiadającymi im hipotezami;
przy remisach wybieramy z góry ustaloną regułę porządku.
Dla $0<\alpha<1$ sprawdzamy kolejno
$p_{(i)}\le\alpha/(m-i+1)$. Odrzucamy hipotezy
$H_{(1)},\ldots,H_{(r)}$ aż do pierwszej niespełnionej
nierówności, po czym zatrzymujemy procedurę.
Jeśli pierwsza nierówność nie zachodzi, nie odrzucamy
żadnej hipotezy.
\end{definicja}

\begin{twierdzenie}[Silna kontrola FWER przez Holma]
\label{thm:holm}
Załóżmy, że każda prawdziwa hipoteza zerowa w rodzinie
$H_1,\ldots,H_m$ ma ważną wartość p, to znaczy
$P(p_i\le u)\le u$ dla każdego $u\in[0,1]$ i
$i\in I_0$. Procedura Holma z definicji~\ref{def:holm}
ma silną kontrolę $\mathrm{FWER}\le\alpha$ bez względu
na zależności między wartościami p. Ponadto zawsze
odrzuca wszystkie hipotezy, które odrzuca Bonferroni
przy tym samym poziomie $\alpha$.
\end{twierdzenie}
\begin{proof}
Jeśli $m_0=0$, to $V=0$ i pierwsza teza jest oczywista.
Niech $m_0\ge1$. Jeżeli Holm odrzuci przynajmniej jedną
prawdziwą hipotezę, niech $j$ będzie jej pierwszą pozycją
w uporządkowanym ciągu. Przed pozycją $j$ znajdują się
wyłącznie hipotezy fałszywe, więc $j\le m-m_0+1$ i
$m-j+1\ge m_0$. Aby Holm odrzucił $H_{(j)}$, musiała
zajść nierówność
\[
 p_{(j)}\le\frac\alpha{m-j+1}\le\frac\alpha{m_0}.
\]
Zatem zdarzenie przynajmniej jednego fałszywego
odrzucenia implikuje, że dla pewnego $i\in I_0$ mamy
$p_i\le\alpha/m_0$. Stosując nierówność sumy do
$m_0$ prawdziwych hipotez, dostajemy $P(V\ge1)\le\alpha$.

Jeśli Bonferroni odrzuca $H_{(j)}$, to
$p_{(j)}\le\alpha/m$. Dla każdego $i\le j$
mamy $p_{(i)}\le p_{(j)}\le\alpha/m\le
\alpha/(m-i+1)$, a zatem Holm nie zatrzyma się
przed pozycją $j$ i również odrzuci $H_{(j)}$.
\end{proof}

\begin{przyklad}[Cztery porównania zaplanowane z góry]
Przypuśćmy, że $m=4$, $\alpha=0{,}05$ i obserwowane
wartości p wynoszą
$0{,}003$, $0{,}018$, $0{,}041$, $0{,}20$.
Bonferroni wymaga $p\le0{,}0125$, więc odrzuca
jedną hipotezę. Holm porównuje kolejno wartości p z
$0{,}0125$, $0{,}016\overline6$, $0{,}025$ i
$0{,}05$; zatrzymuje się na drugiej pozycji i również
odrzuca jedną hipotezę. W innej konfiguracji Holm może
odrzucić więcej hipotez, ale nigdy mniej niż Bonferroni.
\end{przyklad}

\begin{uwaga}[Korekta Šidáka]
Jeżeli \emph{wszystkie} wartości p z prawdziwych hipotez
zerowych są niezależne, to odrzucanie przy
$p_i\le1-(1-\alpha)^{1/m}$ daje silną kontrolę FWER.
Rzeczywiście, dla $m_0\le m$ niezależnych i ważnych
wartości p prawdopodobieństwo braku błędnych odrzuceń
jest co najmniej $(1-t)^{m_0}\ge(1-t)^m=1-\alpha$.
Bez odpowiednich założeń o zależności tej korekty nie
należy stosować jako bezwarunkowej gwarancji.
\end{uwaga}

\section{Benjamini--Hochberg: kontrolowanie odsetka fałszywych odkryć}
\begin{definicja}[Procedura Benjamini--Hochberg (BH)]
\label{def:bh}
Niech $p_{(1)}\le\cdots\le p_{(m)}$ i ustalmy
$q\in(0,1)$. Wyznaczamy
\[
 R=\max\left(\{0\}\cup
 \left\{i\in\{1,\ldots,m\}:
 p_{(i)}\le\frac{iq}{m}\right\}\right).
\]
Jeżeli $R=0$, nie odrzucamy żadnej hipotezy.
Jeżeli $R\ge1$, odrzucamy wszystkie hipotezy z
$p_i\le p_{(R)}$. Liczba tych odrzuceń wynosi dokładnie
$R$: jeśli $R<m$ i $p_{(R+1)}=p_{(R)}$, to również
$p_{(R+1)}\le qR/m<q(R+1)/m$, wbrew maksymalności $R$.
Dla $R=m$ odrzucamy wszystkie $m$ hipotez.
\end{definicja}

\begin{uwaga}[Dlaczego procedura jest krokowa, a nie zbiorem
niezależnych testów?]
Wartości progowe $iq/m$ rosną wraz z indeksem.
Może się zdarzyć, że nierówność nie zachodzi dla pewnego
$i<R$, a jednak zachodzi dla $R$. Procedura BH nie
przerywa badania przy pierwszym niepowodzeniu; wybiera
\emph{największy} indeks spełniający nierówność.
W przeciwieństwie do Holma nie jest procedurą
zatrzymywaną na pierwszym przekroczeniu.
\end{uwaga}

\begin{twierdzenie}[Kontrola FDR przez BH przy niezależności]
\label{thm:bh}
Niech $p_1,\ldots,p_m$ będą niezależnymi zmiennymi
losowymi o wartościach w $[0,1]$. Zakładamy, że
każde $p_i$ odpowiadające prawdziwej hipotezie zerowej
jest ważną wartością p:
$P(p_i\le u)\le u$ dla każdego $u\in[0,1]$.
Wtedy procedura BH na poziomie $q\in(0,1)$ spełnia
\[
 \mathrm{FDR}\le\frac{m_0}{m}q\le q.
\]
\end{twierdzenie}
\begin{proof}
Jeśli $m_0=0$, to FDR wynosi zero. Ustalmy
$i\in I_0$, a następnie pozostawmy wszystkie pozostałe
wartości p bez zmian i zastąpmy $p_i$ przez $0$.
Niech $R_i^0\ge1$ oznacza liczbę odrzuceń w tej
zmodyfikowanej próbie. Z własności krokowej konstrukcji
BH wynika, że na zdarzeniu, w którym oryginalna
procedura odrzuca $H_i$, zachodzi $R=R_i^0$.
Rzeczywiście, ustawienie już odrzuconej wartości p na
zero nie może zmniejszyć maksymalnego indeksu $R$;
nie może też go zwiększyć: pozostałe wartości nie
zmieniły się, a żaden indeks większy niż $R$ nie
spełniał wcześniej warunku progowego.

Warunek odrzucenia $H_i$ pociąga
$p_i\le qR_i^0/m$. Zatem, warunkując względem
wszystkich $p_j$ z $j\ne i$ i używając ich
niezależności od $p_i$, otrzymujemy
\[
 \E\left[\frac{\ind{\{H_i\ \text{odrzucona}\}}}{R\vee1}
 \ \middle|\ (p_j)_{j\ne i}\right]
 \le\frac1{R_i^0}P\left(p_i\le\frac{qR_i^0}{m}
 \ \middle|\ (p_j)_{j\ne i}\right)
 \le\frac q m.
\]
Ostatnia nierówność wynika z ważności $p_i$ i
$1\le R_i^0\le m$. Sumując po $i\in I_0$
i biorąc wartości oczekiwane, dostajemy
$\mathrm{FDR}=\sum_{i\in I_0}
\E[\ind{\{H_i\ \text{odrzucona}\}}/(R\vee1)]
\le m_0q/m$, co dowodzi twierdzenia.
\end{proof}

\begin{przyklad}[Te same cztery testy, inny cel]
Dla wartości p $0{,}003$, $0{,}018$, $0{,}041$, $0{,}20$
i $q=0{,}05$ progi BH to
$0{,}0125$, $0{,}025$, $0{,}0375$, $0{,}05$.
Największym poprawnym indeksem jest $R=2$, więc BH
odrzuca dwie hipotezy, podczas gdy Bonferroni i Holm
odrzucają jedną. Nie oznacza to, że dokładnie
$5\%$ tych dwóch odkryć jest fałszywych; kontrola
dotyczy wartości oczekiwanej FDP w powtarzanym
schemacie badania i opiera się na założeniach
twierdzenia~\ref{thm:bh}.
\end{przyklad}

\section{Zależne testy: co pozostaje prawdziwe?}
Wiele testów może korzystać z tych samych obserwacji,
na przykład z tej samej grupy kontrolnej albo z silnie
skorelowanych pomiarów. Nie wolno wtedy bez komentarza
stosować dowodu BH, który wykorzystywał niezależność
$p_i$ od pozostałych wartości p.

\begin{twierdzenie}[Korekta Benjamini--Yekutieli]\label{thm:by}
Niech $p_1,\ldots,p_m$ będą dowolnie zależnymi
wartościami p, przy czym dla każdego $i\in I_0$
i $u\in[0,1]$ zachodzi $P(p_i\le u)\le u$.
Niech $H_m=\sum_{r=1}^m1/r$. Jeśli w procedurze BH
z definicji~\ref{def:bh} zastąpimy próg $iq/m$ progiem
$iq/(mH_m)$, to otrzymana procedura spełnia
$\mathrm{FDR}\le(m_0/m)q\le q$.
\end{twierdzenie}
\begin{proof}
Niech $c=q/(mH_m)$ i dla ustalonego $i\in I_0$
rozważmy zdarzenia $A_r=\{(r-1)c<p_i\le rc\}$
dla $r=1,\ldots,m$. Jeśli hipoteza $H_i$ jest
odrzucona i $p_i\in A_r$, to z warunku krokowego
$p_i\le cR$ wynika $R\ge r$. Dla $p_i=0$
też zachodzi $R\ge1$ na zdarzeniu odrzucenia;
ponieważ rozkład ważnej wartości p spełnia
$P(p_i=0)=0$, możemy ten przypadek pominąć.
Wobec tego
\[
 \E\frac{\ind{\{H_i\ \text{odrzucona}\}}}{R\vee1}
 \le\sum_{r=1}^m\frac1rP(A_r).
\]
Niech $F_i(t)=P(p_i\le t)$, a $F_i(0)=0$.
Ważność oznacza $F_i(rc)\le rc$. Z przekształcenia
sumy przez części (współczynniki $1/r$ maleją) wynika
\[
 \sum_{r=1}^m\frac{F_i(rc)-F_i((r-1)c)}r
 \le\sum_{r=1}^m\frac{rc-(r-1)c}{r}
 =cH_m=\frac qm.
\]
Dla jasności, nierówność wynika z tożsamości
$\sum_{r=1}^m a_r(F_r-F_{r-1})=
 a_mF_m+\sum_{r=1}^{m-1}(a_r-a_{r+1})F_r$
i oszacowań $F_r\le rc$; suma podstawionych górnych
ograniczeń teleskopuje się do $c\sum_{r=1}^m1/r$.
Sumujemy teraz po $i\in I_0$, otrzymując
$\mathrm{FDR}\le m_0q/m$.
\end{proof}

Korekta BY jest bezpieczna przy dowolnej zależności,
ale może istotnie zmniejszać liczbę odkryć. Dla pewnych
rodzin zależnych testów BH również kontroluje FDR
(np. przy określonych warunkach dodatniej zależności),
lecz nie należy przypisywać tej gwarancji wszystkim
możliwym zależnościom.

\section{Po ANOVA: porównania post hoc i pytania pierwotne}
W rozdziale~\ref{ch:parametric-tests} test ANOVA badał
$H_0:\mu_1=\cdots=\mu_k$. Po jego odrzuceniu często
zadajemy bardziej szczegółowe pytania: czy A różni się od B,
czy A różni się od C, czy B różni się od C? To są
\emph{trzy oddzielne hipotezy}, a ich zbiór powinien zostać
uwzględniony w kontroli błędów.

\begin{center}\small
\begin{tabular}{>{\raggedright\arraybackslash}p{0.27\textwidth}>{\raggedright\arraybackslash}p{0.34\textwidth}>{\raggedright\arraybackslash}p{0.29\textwidth}}
\toprule
\textbf{Sytuacja} & \textbf{Metoda} & \textbf{Dlaczego / ograniczenia}\\
\midrule
Kilka zaplanowanych kontrastów & Bonferroni lub Holm & Silna kontrola FWER przy ważnych wartościach p; nie wymagają niezależności.\\
Wszystkie pary grup w normalnej ANOVA ze wspólną wariancją & Tukey HSD & Wykorzystuje wspólny rozkład maksimum porównań i zwykle jest mniej konserwatywny niż Bonferroni.\\
Wiele grup przeciw jednej kontroli & Dunnett & Korzysta z zależności porównań do wspólnej kontroli.\\
Wiele par przy różnych wariancjach & Games--Howell & Aproksymacja dla wariancji nierównych; należy ocenić liczebności i ogony.\\
Setki testów przesiewowych & BH, ewentualnie BY & Wybór FDR zamiast FWER jest decyzją o tolerowanej skali błędów.\\
\bottomrule
\end{tabular}
\end{center}

Wszystkie te metody wymagają sprawdzenia modelu
prawdopodobieństwa. Nie należy interpretować
„braku istotności pary A--B” jako dowodu równoważności;
podobnie jak w pojedynczym teście konieczne są
odrębne hipotezy równoważności i uzasadniony margines.
Pełną teorię procedur Tukeya, Dunnetta i Gamesa--Howella,
w tym ich statystyki i rozkłady, rozwiniemy w rozdziale
o rozbudowanej ANOVA i ANCOVA.

\section{Praktyczne ryzyka: eksploracja, selekcja i podglądanie danych}
\begin{przyklad}[Wybieranie najlepszego wyniku]
Analityk oblicza pięćdziesiąt korelacji i raportuje
jedynie trzy najmniejsze wartości p, nie ujawniając
pozostałych testów. Odbiorca mógłby błędnie odczytać je
jako trzy z góry zaplanowane porównania. W rzeczywistości
zbiór hipotez jest szerszy; aby ocenić wiarygodność
odkryć, potrzebujemy informacji o całej rodzinie
przeprowadzonych testów. Procedura korekty oparta tylko
na trzech wybranych wynikach nie kontroluje błędów
w pierwotnym schemacie wyboru.
\end{przyklad}

\begin{uwaga}[Sekwencyjne zaglądanie do danych]
Korekty Bonferroniego, Holma i BH rozważane tutaj
dotyczą ustalonej rodziny ważnych wartości p.
Jeśli dla \emph{tej samej} hipotezy wielokrotnie zaglądamy
do narastającej próby i kończymy badanie, gdy zwykłe
$p\le0{,}05$, to końcowa wartość p nie musi już być
ważna przy tej regule zatrzymania. Potrzebne są osobno
zaplanowane metody sekwencyjne: wydatkowanie poziomu
istotności, granice grupowo-sekwencyjne albo testy ważne
w czasie. Samo dodanie korekty na liczbę raportowanych
hipotez nie naprawia takiego problemu.
\end{uwaga}

\section{Dobór procedury w praktyce}
\begin{center}\small
\begin{tabular}{>{\raggedright\arraybackslash}p{0.31\textwidth}>{\raggedright\arraybackslash}p{0.28\textwidth}>{\raggedright\arraybackslash}p{0.29\textwidth}}
\toprule
\textbf{Pytanie} & \textbf{Zwykle rozważamy} & \textbf{Na co uważać?}\\
\midrule
Czy choć jeden fałszywy alarm byłby kosztowny? & FWER: Holm, Bonferroni & Silna kontrola całej rodziny.\\
Czy wykonujemy liczne badania przesiewowe? & FDR: BH & Niezależność lub odpowiednie warunki zależności.\\
Czy zależności p-value są nieznane? & Holm dla FWER, BY dla FDR & BY może mieć małą moc.\\
Czy porównujemy wszystkie pary po ANOVA? & Tukey lub Games--Howell & Wspólna wariancja lub jej brak, liczebności.\\
Czy wszystkie grupy porównujemy z kontrolą? & Dunnett & Wspólny model i współzależność kontrastów.\\
Czy obserwowaliśmy dane wielokrotnie w czasie? & Procedura sekwencyjna & Predefiniowany plan, ważność przy zatrzymaniu.\\
Czy raportujemy wiele oszacowań? & Jednoczesne przedziały & Własności pokrycia całej rodziny.\\
\bottomrule
\end{tabular}
\end{center}

\begin{lstlisting}[language=Python,caption={Kompletny program w Pythonie: Bonferroni, Holm, BH i BY bez zewnętrznych pakietów.}]
from math import fsum

pvals = [0.003, 0.018, 0.041, 0.20]
alpha = 0.05
m = len(pvals)
order = sorted(range(m), key=lambda i: pvals[i])

bonf = [p <= alpha/m for p in pvals]
holm = [False] * m
for rank, i in enumerate(order):
    if pvals[i] > alpha/(m-rank):
        break
    holm[i] = True

def step_up(pvals, q):
    m = len(pvals)
    order = sorted(range(m), key=lambda i: pvals[i])
    r = max((j for j in range(1, m+1)
             if pvals[order[j-1]] <= q*j/m), default=0)
    return [p <= pvals[order[r-1]] for p in pvals] if r else [False]*m

harmonic = fsum(1/j for j in range(1, m+1))
bh = step_up(pvals, alpha)
by = step_up(pvals, alpha/harmonic)
print('Bonferroni:', bonf)
print('Holm:', holm)
print('BH:', bh)
print('BY:', by)
\end{lstlisting}

\begin{lstlisting}[language=Julia,caption={Obliczenie krokowych korekt Holma i BH w Julii bez pakietów zewnętrznych.}]
pvals = [0.003, 0.018, 0.041, 0.20]
alpha = 0.05
m = length(pvals)
idx = sortperm(pvals)

holm = falses(m)
for j in 1:m
    i = idx[j]
    pvals[i] > alpha/(m-j+1) && break
    holm[i] = true
end

threshold = [alpha*j/m for j in 1:m]
valid = [j for j in 1:m if pvals[idx[j]] <= threshold[j]]
r = isempty(valid) ? 0 : maximum(valid)
bh = falses(m)
if r > 0
    bh .= pvals .<= pvals[idx[r]]
end
println((holm=holm, bh=bh))
\end{lstlisting}

\paragraph{Jak czytać wyniki programu?}
Obie procedury uruchomione na tych samych danych mogą
zwrócić różne decyzje, mimo że każda jest poprawna
\emph{względem swojego celu i założeń}. Nie należy
wybierać między FWER a FDR dopiero na podstawie
tego, która procedura odrzuciła więcej hipotez.

\section{Zadanie podsumowujące: rekomendacje dla analityka}
\begin{zadanie}[Cztery lokalizacje, trzy rodzaje pytań]
Firma porównuje cztery lokalizacje pod kątem wskaźników
awaryjności. Rozważa trzy cele: (a) potwierdzić, że
\emph{każda} wskazana lokalizacja spełnia normę;
(b) znaleźć możliwie wiele lokalizacji wymagających
interwencji; (c) porównać każdą lokalizację z jedną
wspólną lokalizacją kontrolną. Jakie metody rozważyć?
\par\emph{Rozwiązanie.} W (a) wybieramy procedurę
kontrolującą rodzinny błąd zgodnie z kierunkiem hipotez,
np. Holm, oraz odpowiednio jednostronne przedziały
jednoczesne. W (b), jeśli akceptujemy kontrolę
oczekiwanego udziału błędnych sygnałów i rozumiemy
koszty pomyłek, możemy rozważyć BH lub BY zależnie od
struktury zależności. W (c) porównania do wspólnej
kontroli sugerują procedurę Dunnetta \emph{w odpowiednim
modelu normalnym}; przy niezależnych prawidłowych
wartościach p albo bez szczególnego modelu korelacji
można użyć Holma. W każdym przypadku od wyboru
procedury ważniejsza jest poprawność jednostki analizy,
założeń modelu i interpretacji wyniku.
\end{zadanie}



% ================================================================
% ROZDZIAL 13 -- REGRESJA LINIOWA
% ================================================================
